\documentclass[12pt,a4paper]{article}

% ============================================================
% Paquetes
% ============================================================
\usepackage[utf8]{inputenc}
\usepackage[T1]{fontenc}
\usepackage[spanish]{babel}
\usepackage{amsmath,amssymb,amsfonts}
\usepackage{graphicx}
\usepackage{xcolor}
\usepackage{hyperref}
\usepackage{geometry}
\usepackage{fancyhdr}
\usepackage{titlesec}
\usepackage{float}
\usepackage{caption}
\usepackage{booktabs}
\usepackage{longtable}
\usepackage{array}
\usepackage{verbatim}

\geometry{margin=2.5cm}
\setlength{\headheight}{13.6pt}

% ============================================================
% Configuracion de encabezado/pie
% ============================================================
\pagestyle{fancy}
\fancyhf{}
\fancyhead[L]{\small 3Dosim --- Modulo 1}
\fancyhead[R]{\small Segmentacion, Registro y Fusion}
\fancyfoot[C]{\thepage}
\renewcommand{\headrulewidth}{0.4pt}

% ============================================================
% Macro para notacion: define una variable con descripcion
% Uso: \vardef{$A_{vox}$}{Concentracion de actividad en Bq/mL}
% ============================================================
\newcommand{\vardef}[2]{\noindent\textbf{#1}: #2\\}

\setlength{\parskip}{0.5em}

% ============================================================
% Portada
% ============================================================
\title{\textbf{Modulo 1: Segmentacion, Registro y Fusion de Imagenes}\\[0.5cm]
       \large Pipeline de Preparacion para Dosimetria Hepatica con $^{90}$Y\\[0.3cm]
       \normalsize Documentacion Tecnica --- Revision v4.0}
\author{Proyecto 3Dosim --- Dosimetria Hepatica con $^{90}$Y}
\date{Julio 2026}

\begin{document}

\maketitle
\thispagestyle{empty}
\newpage

% ============================================================
% GLOSARIO DE ACRONIMOS
% ============================================================
\section*{Glosario de Acronimos}
\addcontentsline{toc}{section}{Glosario de Acronimos}

\begin{description}
    \item[3Dosim] Pipeline de dosimetria interna 3D para radioembolizacion hepatica.
    \item[A$_{\text{total}}$] Actividad total integrada sobre todo el volumen PET [Bq].
    \item[ASGD] Adaptive Stochastic Gradient Descent --- optimizador de Elastix.
    \item[Bq] Becquerel --- unidad de actividad radiactiva (1 desintegracion/s).
    \item[Bq/mL] Concentracion de actividad por mililitro de tejido.
    \item[CT] Computed Tomography --- imagen anatomica en unidades Hounsfield (HU).
    \item[DICOM] Digital Imaging and Communications in Medicine --- estandar
          de almacenamiento y transmision de imagenes medicas.
    \item[DOF] Degrees of Freedom --- grados de libertad de una transformacion.
    \item[FMESH4] Mesh tally multiplicador de MCNP --- registra energia
          depositada por voxel.
    \item[FOV] Field of View --- campo de vision de la imagen.
    \item[FU] Fraction of Uptake --- fraccion de actividad captada por un organo.
    \item[GBq] Gigabecquerel ($10^9$ Bq). Dosis tipica de $^{90}$Y: 1--5 GBq.
    \item[GDPR] General Data Protection Regulation (Union Europea).
    \item[HIPAA] Health Insurance Portability and Accountability Act (EEUU).
    \item[HU] Hounsfield Unit --- escala de densidad radiologica:
          agua $= 0$ HU, aire $= -1000$ HU, hueso $> 400$ HU.
    \item[ICRP-110] Publicacion 110 de la Comision Internacional de
          Proteccion Radiologica --- phantoms computacionales de referencia.
    \item[ICRU-44] Reporte 44 de la Comision Internacional de Unidades y
          Medidas Radiologicas --- composiciones de tejidos.
    \item[MAE] Mean Absolute Error.
    \item[MCNP] Monte Carlo N-Particle --- codigo de simulacion de transporte
          de radiacion (Los Alamos National Laboratory).
    \item[MCTAL] Archivo de salida de MCNP con resultados de tallies.
    \item[mCi] Milicurie ($1$ mCi $= 37$ MBq $= 3.7 \times 10^7$ Bq). Unidad historica.
    \item[MI] Mutual Information --- metrica de similitud entre imagenes.
    \item[MIRD] Medical Internal Radiation Dose --- modelo analitico de dosimetria.
    \item[MSE] Mean Squared Error --- metrica de diferencia cuadratica media.
    \item[NCC] Normalized Cross-Correlation --- metrica de similitud entre imagenes.
    \item[NIfTI] Neuroimaging Informatics Technology Initiative --- formato
          de archivo volumetrico (.nii), usado por MCNP.
    \item[NPS] Number of Particle Histories --- numero de historias en MCNP.
    \item[NRRD] Nearly Raw Raster Data --- formato volumetrico (.nrrd),
          nativo de 3D Slicer.
    \item[PET] Positron Emission Tomography --- imagen funcional de actividad.
    \item[PHI] Protected Health Information --- datos de salud protegidos (HIPAA).
    \item[RMSE] Root Mean Squared Error.
    \item[ROI] Region of Interest.
    \item[SF] Shunt Fraction --- fraccion de microesferas que escapan al pulmon.
    \item[SUV] Standardized Uptake Value --- medida semicuantitativa PET.
    \item[T/N] Tumor-to-Normal ratio --- razon de captacion PET tumor/higado sano.
    \item[TS] TotalSegmentator --- segmentador anatomico por deep learning
          (104 clases a partir de CT).
    \item[nnU-Net] Arquitectura de red neuronal convolucional usada por TS.
\end{description}

\newpage

% ============================================================
% Tabla de contenidos
% ============================================================
\tableofcontents
\newpage

% ============================================================
% 1. PIPELINE COMPLETO
% ============================================================
\section{Pipeline Completo --- Los 16 Pasos}

\subsection{Vision General}

La dosimetria interna con $^{90}$Y requiere conocer exactamente donde esta cada tejido (higado,
tumor, pulmon, hueso) y cuanta actividad hay en cada punto (PET). El Modulo 1 toma imagenes
DICOM crudas --- un CT anatomico y un PET funcional --- y las procesa a traves de 16 pasos
automatizados para producir una \textbf{labelmap dosimetrica}: un volumen 3D donde cada voxel
tiene un valor entero que identifica su tipo de tejido segun la tabla de materiales
ICRP-110/ICRU-44. Esta labelmap es la entrada del Modulo 2 (generacion de entrada MCNP).

\subsection*{Acronimos usados en esta seccion}
\begin{tabular}{ll}
\toprule
\textbf{Acronimo} & \textbf{Significado} \\
\midrule
CT & Computed Tomography \\
DICOM & Digital Imaging and Communications in Medicine \\
HU & Hounsfield Unit \\
MCNP & Monte Carlo N-Particle \\
NIfTI & Neuroimaging Informatics Technology Initiative \\
NRRD & Nearly Raw Raster Data \\
PET & Positron Emission Tomography \\
TS & TotalSegmentator \\
\bottomrule
\end{tabular}

\subsection{Diagrama de Flujo General}

\begin{verbatim}
+-----------------------------------------------------------------+
|                         MODULO 1                                |
|              Segmentacion + Registro + Fusion                    |
|                                                                  |
|  DICOM ---> Calibracion ---> Camilla ---> Registro ---> Fusion  |
|   CT       PET            removal     PET->CT      CT+PET       |
|    |          |              |            |           |          |
|    +----------+--------------+------------+-----------+          |
|                           |                                      |
|                           v                                      |
|  Anonimizacion ---> TotalSegmentator ---> Validacion IA          |
|       |                  |                        |              |
|       +------------------+------------------------+              |
|                           |                                      |
|                           v                                      |
|  Validacion Medica ---> Creacion Tumor ---> Validacion Tumor     |
|                           |                                      |
|                           v                                      |
|  Higado Sano ---> Body seg. ---> Labelmap NIfTI+NRRD             |
|                                                                  |
|                              v                                   |
|                    Labelmap Dosimetrica                          |
|                    (entrada Modulo 2)                            |
+-----------------------------------------------------------------+
\end{verbatim}

\subsection{Secuencia de Pasos}

Cada paso tiene un archivo fuente, un proposito y una salida especifica.
A continuacion se listan los 16 pasos con referencias a las secciones detalladas.

\begin{enumerate}
    \item \textbf{Check Slicer} (Seccion \ref{sec:check}): Verifica entorno, modulos y rutas.
    \item \textbf{Carga DICOM} (Seccion \ref{sec:carga}): Carga CT y PET desde directorios DICOM.
    \item \textbf{Calibracion PET} (Seccion \ref{sec:calibracion}): Convierte DICOM raw a Bq/mL.
    \item \textbf{Eliminar Camilla} (Seccion \ref{sec:camilla}): Remueve mesa y aire del CT.
    \item \textbf{Registro PET$\to$CT} (Seccion \ref{sec:registro}): Re-muestrea PET a grilla CT.
    \item \textbf{Fusion CT+PET} (Seccion \ref{sec:fusion}): Layout 4-up + dialogo de actividad.
    \item \textbf{Anonimizacion} (Seccion \ref{sec:anonimizacion}): Elimina metadatos PHI.
    \item \textbf{Exportar metadatos} (Seccion \ref{sec:anonimizacion}): Metadatos DICOM a CSV.
    \item \textbf{TotalSegmentator} (Seccion \ref{sec:ts}): Segmentacion 104 clases (nnU-Net).
    \item \textbf{Validacion IA} (Seccion \ref{sec:ai}): AI Supervisor verifica automaticamente.
    \item \textbf{Validacion Medica} (Seccion \ref{sec:validacion}): Aprobacion obligatoria.
    \item \textbf{Creacion de Tumor} (Seccion \ref{sec:tumor}): 4 modos de lesion hepatica.
    \item \textbf{Validacion de Tumor} (Seccion \ref{sec:validacion_tumor}): Aprobacion medica.
    \item \textbf{Higado Sano} (Seccion \ref{sec:tumor}): Higado completo menos tumor.
    \item \textbf{Body Segmentation} (Seccion \ref{sec:body}): Contorno externo con TS.
    \item \textbf{Exportar Labelmap} (Seccion \ref{sec:labelmap}): NIfTI + NRRD con indices.
\end{enumerate}

\subsection{Modos de Ejecucion}

\begin{table}[H]
\centering
\small
\begin{tabular}{p{2.8cm} p{7.5cm} p{2.2cm}}
\toprule
\textbf{Modo} & \textbf{Comando} & \textbf{Pasos} \\
\midrule
Completo & \texttt{Slicer.exe --python-script main.py --data-dir ''Paciente\_2''} & 16 \\
Hasta fusion & \texttt{--stop-after-fusion} & 1--6 \\
Antes de segmentar & \texttt{--stop-before-segment} & 1--7 \\
Reset & \texttt{--reset} & 16 (desde cero) \\
Debug & \texttt{--debug} & 16 (pausa tras cada paso) \\
\bottomrule
\end{tabular}
\end{table}

\subsection{Estructura de Archivos de Salida}

\begin{verbatim}
output_dir/
+-- anon/
|   +-- ANON_CT.nrrd       (CT anonimizado)
|   +-- ANON_PET.nii       (PET anonimizado)
+-- checkpoints/
|   +-- pipeline_checkpoint.json
+-- labelmaps/
|   +-- labelmap.nii       (NIfTI para Modulo 2)
|   +-- labelmap.nrrd      (NRRD para Slicer)
+-- scenes/
|   +-- 3Dosim_scene.mrb   (escena Slicer)
+-- screenshots/
+-- logs/
|   +-- pipeline_mod1.log
+-- summary.txt
+-- pipeline_results.json  (historial de ejecuciones)
\end{verbatim}

\newpage

% ============================================================
% 1b. CHECK SLICER
% ============================================================
\section{Verificacion de Entorno (Check Slicer)}
\label{sec:check}

\subsection*{Acronimos usados en esta seccion}
\begin{tabular}{ll}
\toprule
\textbf{Acronimo} & \textbf{Significado} \\
\midrule
CLI & Command-Line Interface \\
TS & TotalSegmentator \\
\bottomrule
\end{tabular}

La ejecucion del pipeline comienza verificando que 3D Slicer este funcionando correctamente, que
los modulos necesarios esten disponibles (BRAINSResample, TotalSegmentator) y que los directorios
de entrada/salida existan. Si el flag \texttt{--reset} esta activo, se eliminan todos los
checkpoints previos para forzar una ejecucion completa desde cero.

\subsection{Verificaciones}

\begin{enumerate}
    \item Comprobar que \texttt{slicer} este importado correctamente.
    \item Verificar existencia de modulos criticales: BRAINSResample, TotalSegmentator.
    \item Validar directorios de entrada: $\mathcal{D}_{\text{CT}}$, $\mathcal{D}_{\text{PET}}$.
    \item Si \texttt{--reset}: eliminar \texttt{pipeline\_checkpoint.json} y escenas \texttt{.mrb}.
\end{enumerate}

\subsection{Control de Calidad (AI Supervisor)}

\begin{table}[H]
\centering
\begin{tabular}{llc}
\toprule
\textbf{Verificacion} & \textbf{Metrica} & \textbf{Umbral} \\
\midrule
Slicer disponible & \texttt{slicer.\_\_version\_\_} existe & Critico \\
Modulo BRAINSResample & Cargado en \texttt{slicer.modules} & Critico \\
Modulo TotalSegmentator & Cargado en \texttt{slicer.modules} & Critico \\
Directorios de entrada & \texttt{os.path.exists($\mathcal{D}$)} & Critico \\
\bottomrule
\end{tabular}
\end{table}

\newpage

% ============================================================
% 2. CARGA DICOM
% ============================================================
\section{Carga DICOM CT+PET}
\label{sec:carga}

\subsection*{Acronimos usados en esta seccion}
\begin{tabular}{ll}
\toprule
\textbf{Acronimo} & \textbf{Significado} \\
\midrule
CT & Computed Tomography \\
DB & Database (base de datos) \\
DICOM & Digital Imaging and Communications in Medicine \\
FOV & Field of View \\
HU & Hounsfield Unit \\
PET & Positron Emission Tomography \\
UID & Unique Identifier \\
\bottomrule
\end{tabular}

La dosimetria con $^{90}$Y requiere dos tipos de imagenes: un CT anatomico (alta resolucion,
$512\times512$ voxels, $0.976$ mm de espaciado) y un PET funcional (menor resolucion,
$200\times200$ voxels, $4.07$ mm de espaciado). Ambas provienen de equipos diferentes y se
almacenan en archivos DICOM separados. Este paso las indexa en una base de datos temporal de
Slicer, construye los volumenes 3D a partir de los slices individuales y renombra los nodos con
nombres canonicos para que el resto del pipeline pueda referenciarlos consistentemente incluso
tras restaurar un checkpoint.

\subsection{Algoritmo Detallado}

Sea $\mathcal{D}_{\text{CT}}$ el directorio con archivos DICOM del CT
y $\mathcal{D}_{\text{PET}}$ el directorio con archivos DICOM del PET.
El proceso de carga sigue estos pasos:

\begin{enumerate}
    \item Verificar existencia de $\mathcal{D}_{\text{CT}}$ y $\mathcal{D}_{\text{PET}}$.
    \item Abrir base de datos DICOM temporal:
          \texttt{DICOMUtils.openTemporaryDatabase()}.
    \item Indexar cada directorio:
          \texttt{DICOMUtils.importDicom($\mathcal{D}$)}.
          Esto parsea los archivos .dcm y extrae metadatos (PatientID,
          StudyDate, Modality, SeriesUID, etc.).
    \item Obtener UIDs de series disponibles:
          \texttt{DICOMUtils.allSeriesUIDsInDatabase()}.
    \item Cargar volumenes:
          \texttt{DICOMUtils.loadSeriesByUID(uids)}.
          Slicer construye arrays 3D a partir de los slices individuales.
    \item Identificar CT vs PET por nombre de nodo.
          Si el nombre contiene \texttt{``CT''} $\to$ CT;
          si contiene \texttt{``PET''} o \texttt{``PT''} o \texttt{``NM''} $\to$ PET.
    \item Renombrar nodos a nombres canonicos \texttt{``CT''} y \texttt{``PET''}
          para consistencia entre checkpoints.
\end{enumerate}

\subsection{Parametros Tipicos de las Imagenes}

\begin{table}[H]
\centering
\begin{tabular}{l c c}
\toprule
\textbf{Parametro} & \textbf{CT tipico} & \textbf{PET tipico} \\
\midrule
Modalidad DICOM & CT & PT \\
Dimensiones ($N_x \times N_y \times N_z$) & $512 \times 512 \times N_s$ & $200 \times 200 \times N_s$ \\
Espaciado ($s_x, s_y, s_z$) [mm] & $0.976 \times 0.976 \times 3.0$ & $4.07 \times 4.07 \times 2.0$ \\
Tipo de dato & int16 (enteros con signo) & float32 \\
Rango de valores & $-1024$ a $+3000$ HU & $0$ a $10^5$ Bq/mL \\
Bits por pixel & 16 & 16 o 32 \\
Campo de vision (FOV) & 500 mm & 814 mm \\
\bottomrule
\end{tabular}
\end{table}

\subsection{Notacion de Variables}

\begin{table}[H]
\centering
\begin{tabular}{lll}
\toprule
\textbf{Variable} & \textbf{Descripcion} & \textbf{Unidad} \\
\midrule
$N_s$ & Numero de slices (tipicamente 100--600) & --- \\
$s_x, s_y, s_z$ & Espaciado entre voxeles & mm \\
$I_{\text{CT}}(x,y,z)$ & Volumen CT en HU & HU \\
$I_{\text{PET}}(x,y,z)$ & Volumen PET en unidades raw DICOM & raw \\
\bottomrule
\end{tabular}
\end{table}

\subsection{Formato DICOM}

Cada archivo DICOM contiene un slice individual con metadatos organizados en tags:

\begin{table}[H]
\centering
\begin{tabular}{lll}
\toprule
\textbf{Categoria} & \textbf{Tag DICOM} & \textbf{Campo} \\
\midrule
Metadatos & (0010,0010) & PatientName \\
 & (0010,0020) & PatientID \\
 & (0008,0020) & StudyDate \\
 & (0008,0060) & Modality \\
Geometria & (0028,0030) & PixelSpacing ($s_x, s_y$) \\
 & (0018,0050) & SliceThickness ($s_z$) \\
 & (0020,0032) & ImagePositionPatient (origen) \\
 & (0020,0037) & ImageOrientationPatient (cosenos) \\
Calibracion & (0028,1052) & RescaleSlope ($m_k$) \\
 & (0028,1053) & RescaleIntercept ($b_k$) \\
 & (0028,1054) & RescaleType \\
\bottomrule
\end{tabular}
\end{table}

\subsection{Control de Calidad (AI Supervisor)}

\begin{table}[H]
\centering
\begin{tabular}{llc}
\toprule
\textbf{Verificacion} & \textbf{Metrica} & \textbf{Umbral} \\
\midrule
CT detectado & Nodo con ``CT'' en nombre existe & Critico \\
PET detectado & Nodo con ``PET''/``PT''/``NM'' existe & Critico \\
Dimensiones CT & $N_x = 512,\ N_y = 512$ & Warning \\
Dimensiones PET & $N_x \geq 100,\ N_y \geq 100$ & Warning \\
Escena guardada & Escritura exitosa de \texttt{.mrb} & Warning \\
\bottomrule
\end{tabular}
\end{table}

\newpage

% ============================================================
% 3. CALIBRACION PET
% ============================================================
\section{Calibracion PET a Bq/mL}
\label{sec:calibracion}

\subsection*{Acronimos usados en esta seccion}
\begin{tabular}{ll}
\toprule
\textbf{Acronimo} & \textbf{Significado} \\
\midrule
Bq/mL & Becquerel por mililitro --- concentracion de actividad \\
DICOM & Digital Imaging and Communications in Medicine \\
GBq & Gigabecquerel ($10^9$ Bq) \\
mCi & Milicurie ($3.7 \times 10^7$ Bq) \\
PET & Positron Emission Tomography \\
SUV & Standardized Uptake Value \\
\bottomrule
\end{tabular}

El PET DICOM almacena valores crudos (cuentas detectadas) que no son directamente utilizables
para dosimetria. Cada fabricante de escaner PET aplica factores de calibracion diferentes, y
estos factores pueden variar incluso entre slices del mismo estudio. Este paso lee cada slice
DICOM individualmente, extrae los factores \texttt{RescaleSlope} y \texttt{RescaleIntercept}
especificos de cada uno y convierte los valores raw a concentracion de actividad radiactiva en
Bq/mL. Es una replica del algoritmo MATLAB \texttt{f\_Rescale\_Bq.m}.

\subsection{Algoritmo de Calibracion}

Sea $P_{\text{raw}}^{(k)}$ la matriz 2D de valores crudos del slice $k$,
con dimensiones $N_x \times N_y$. La calibracion se aplica voxel a voxel:

\begin{enumerate}
    \item Listar archivos DICOM en $\mathcal{D}_{\text{PET}}$, ordenados por nombre.
    \item Leer cada slice $k$ con \texttt{pydicom.dcmread()}.
    \item Extraer factores del slice $k$: $m_k$, $b_k$, $\tau_k$.
    \item Si $\tau_k = \texttt{``BQML''}$, aplicar la ecuacion de calibracion.
\end{enumerate}

\begin{equation}\label{eq:calibracion}
    A_{\text{vox}}^{(k)}(i,j) = P_{\text{raw}}^{(k)}(i,j) \cdot m_k + b_k
\end{equation}
\vardef{$A_{\text{vox}}^{(k)}(i,j)$}{Concentracion de actividad en Bq/mL en el voxel $(i,j)$ del slice $k$}
\vardef{$P_{\text{raw}}^{(k)}(i,j)$}{Valor crudo DICOM del voxel $(i,j)$ del slice $k$}
\vardef{$m_k$}{RescaleSlope del slice $k$}
\vardef{$b_k$}{RescaleIntercept del slice $k$}

\begin{table}[H]
\centering
\begin{tabular}{lll}
\toprule
\textbf{Variable} & \textbf{Descripcion} & \textbf{Unidad} \\
\midrule
$A_{\text{vox}}^{(k)}(i,j)$ & Concentracion de actividad calibrada & Bq/mL \\
$P_{\text{raw}}^{(k)}(i,j)$ & Valor crudo DICOM & adimensional \\
$m_k$ & RescaleSlope del slice $k$ & (Bq/mL)/unidad \\
$b_k$ & RescaleIntercept del slice $k$ & Bq/mL \\
$\tau_k$ & RescaleType del slice $k$ & --- \\
$i, j$ & Indices de columna y fila & --- \\
$k$ & Indice de slice & --- \\
\bottomrule
\end{tabular}
\end{table}

\subsection{Conversion a Becquerels por Voxel}

De Bq/mL a Bq por voxel:

\begin{equation}\label{eq:bq_por_voxel}
    A_{\text{Bq}}^{(k)}(i,j) = A_{\text{vox}}^{(k)}(i,j) \cdot V_{\text{voxel}}
\end{equation}
\vardef{$A_{\text{Bq}}^{(k)}(i,j)$}{Actividad total en el voxel $(i,j,k)$}
\vardef{$V_{\text{voxel}}$}{Volumen de un voxel PET en mL}

donde el volumen del voxel se calcula a partir del espaciado DICOM:

\begin{equation}\label{eq:vol_voxel}
    V_{\text{voxel}} = \frac{s_x^{\text{PET}}}{10} \cdot
    \frac{s_y^{\text{PET}}}{10} \cdot
    \frac{s_z^{\text{PET}}}{10} \quad [\text{cm}^3 = \text{mL}]
\end{equation}
\vardef{$s_x^{\text{PET}}, s_y^{\text{PET}}$}{PixelSpacing del PET en mm}
\vardef{$s_z^{\text{PET}}$}{SliceThickness del PET en mm}

\subsection{Actividad Total}

La actividad total en el volumen PET es la suma sobre todos los voxeles:

\begin{equation}\label{eq:act_total}
    A_{\text{total}} = \sum_{k=1}^{N_z} \sum_{i=1}^{N_y} \sum_{j=1}^{N_x}
    A_{\text{Bq}}^{(k)}(i,j) \quad [\text{Bq}]
\end{equation}
\vardef{$A_{\text{total}}$}{Actividad total integrada}

Las conversiones a otras unidades:

\begin{align}
    A_{\text{GBq}} &= A_{\text{total}} / 10^9 \quad [\text{GBq}] \\
    A_{\text{mCi}} &= A_{\text{total}} / (3.7 \times 10^7) \quad [\text{mCi}]
\end{align}

donde $1$ mCi $= 37$ MBq $= 3.7 \times 10^7$ Bq.

\subsection{Parametros Tipicos}

\begin{table}[H]
\centering
\begin{tabular}{l c l}
\toprule
\textbf{Parametro} & \textbf{Valor tipico} & \textbf{Descripcion} \\
\midrule
RescaleSlope ($m_k$) & 1.0--2.5 & Factor de pendiente por slice \\
RescaleIntercept ($b_k$) & 0.0 & Intercepto (tipicamente 0) \\
RescaleType ($\tau_k$) & \texttt{``BQML''} & Unidad: Bq/mL \\
Voxel PET ($V_{\text{voxel}}$) & 0.042 mL & Para $4.07 \times 4.07 \times 2.5$ mm \\
Actividad administrada ($^{90}$Y) & 1--5 GBq & Rango tipico clinico \\
Actividad en imagen & 0.5--3 GBq & Depende del tiempo post-administracion \\
Numero de slices & 100--400 & Segun la duracion de la adquisicion \\
\bottomrule
\end{tabular}
\end{table}

\subsection{Conservacion de la Actividad Post-Registro}

Cuando el PET se re-muestrea a la grilla del CT (Paso 5), la interpolacion altera la suma total.
Para compensar se aplica un factor de correccion:

\begin{equation}\label{eq:conservacion}
    \text{PET}_{\text{corregido}}(\mathbf{r}) =
    \text{PET}_{\text{interp}}(\mathbf{r}) \cdot
    \frac{A_{\text{original}}}{A_{\text{interpolada}}}
\end{equation}
\vardef{$\mathbf{r}$}{Coordenada voxel $(x,y,z)$ en la grilla CT}
\vardef{$\text{PET}_{\text{interp}}(\mathbf{r})$}{PET interpolado en Bq/mL}
\vardef{$A_{\text{original}}$}{Actividad total antes del resample [Bq]}
\vardef{$A_{\text{interpolada}}$}{Actividad total del PET interpolado [Bq]}

\subsection{Advertencia Critica: Slicer Pierde la Calibracion por Slice}

El nodo de Slicer \textbf{pierde} los factores RescaleSlope/RescaleIntercept por slice: Slicer
aplica factores globales a todo el volumen, no por archivo. Por esta razon el pipeline lee la
actividad desde DICOM raw con \texttt{pydicom}, no desde el nodo Slicer.

\subsection{Control de Calidad (AI Supervisor)}

\begin{table}[H]
\centering
\begin{tabular}{llc}
\toprule
\textbf{Verificacion} & \textbf{Metrica} & \textbf{Umbral} \\
\midrule
Actividad total & $A_{\text{total}} > 0$ & Critico \\
RescaleType & $\tau_k = \text{``BQML''}$ & Warning \\
Rango de actividad & $0.1 \leq A_{\text{total}} \leq 50$ GBq & Warning \\
Slices validos & $N_{\text{slices}} > 0$ & Critico \\
\bottomrule
\end{tabular}
\end{table}

\newpage

% ============================================================
% 4. ELIMINACION DE CAMILLA
% ============================================================
\section{Eliminacion de Camilla y Aire}
\label{sec:camilla}

\subsection*{Acronimos usados en esta seccion}
\begin{tabular}{ll}
\toprule
\textbf{Acronimo} & \textbf{Significado} \\
\midrule
CT & Computed Tomography \\
FOV & Field of View \\
HU & Hounsfield Unit \\
\bottomrule
\end{tabular}

El CT incluye la camilla (mesa de exploracion) y el aire alrededor del paciente, que no deben
formar parte del modelo dosimetrico. Este paso separa el cuerpo del paciente del fondo mediante
un umbral de Hounsfield Units, seguido de operaciones morfologicas y de conectividad para obtener
una mascara corporal limpia. El CT original se conserva intacto para TotalSegmentator, que
necesita el FOV completo; se crea un nuevo nodo \texttt{CT\_sin\_camilla} con la mascara aplicada.

\subsection{Algoritmo Detallado}

Sea $I_{\text{CT}}(x,y,z)$ el volumen CT original en HU. El proceso genera una mascara corporal
$M(x,y,z)$ binaria ($1$ = cuerpo, $0$ = fondo).

\subsubsection{Paso 1: Umbralizacion}

Se aplica un umbral sobre los valores HU para separar tejido del aire:

\begin{equation}\label{eq:umbral}
    M_1(x,y,z) = \begin{cases}
        1 & \text{si } I_{\text{CT}}(x,y,z) > -200 \\
        0 & \text{en caso contrario}
    \end{cases}
\end{equation}
\vardef{$M_1$}{Mascara binaria inicial (tejido vs.\ aire)}
\vardef{$I_{\text{CT}}$}{Volumen CT original [HU]}

El valor $-200$ HU corresponde al limite inferior del tejido blando.
Solo el aire puro ($<-900$ HU) queda excluido.

\subsubsection{Paso 2: Cierre Morfologico}

Operacion morfologica para rellenar huecos internos (vasos, bronquios) dentro del cuerpo:

\begin{equation}\label{eq:cierre}
    M_2 = (M_1 \oplus B) \ominus B
\end{equation}
\vardef{$M_2$}{Mascara tras cierre morfologico}
\vardef{$\oplus$}{Operador de dilatacion morfologica}
\vardef{$\ominus$}{Operador de erosion morfologica}
\vardef{$B$}{Elemento estructurante cubico $3\times3\times3$}

Parametros: elemento estructurante $B$ de $3 \times 3 \times 3$ voxeles, 3 iteraciones.
Implementacion: \texttt{scipy.ndimage.binary\_closing}.

\subsubsection{Paso 3: Componente Conexa Mayoritaria}

En cada slice axial $z$, se etiquetan las componentes conectadas (conectividad 8-vecina en 2D)
y se conserva solo la mas grande:

\begin{equation}\label{eq:comp_conexa}
    M_3(z,:,:) = \arg\max_{c \in \mathcal{C}_z} |c|
\end{equation}
\vardef{$\mathcal{C}_z$}{Conjunto de componentes conectadas en el slice $z$ de $M_2$}
\vardef{$|c|$}{Numero de voxeles de la componente $c$}
\vardef{$\arg\max$}{Selecciona la componente con mayor $|c|$}

\subsubsection{Paso 4: Eliminacion de Camilla}

En cada slice axial se identifica la fila inferior del cuerpo y se elimina todo lo que esta
debajo (la camilla). Tambien se recortan 2--3 pixeles de los bordes laterales para eliminar
artefactos.

\subsubsection{Paso 5: Aplicacion de Mascara}

El CT enmascarado se construye fijando a $-1024$ HU (aire) los voxeles fuera del cuerpo:

\begin{equation}\label{eq:ct_masked}
    I_{\text{enmasc}}(x,y,z) = \begin{cases}
        I_{\text{CT}}(x,y,z) & \text{si } M(x,y,z) = 1 \\
        -1024 & \text{en caso contrario}
    \end{cases}
\end{equation}
\vardef{$I_{\text{enmasc}}$}{Volumen CT enmascarado [HU]}
\vardef{$M$}{Mascara corporal final (binaria)}

\subsection{Parametros}

\begin{table}[H]
\centering
\begin{tabular}{l c l}
\toprule
\textbf{Parametro} & \textbf{Valor} & \textbf{Descripcion} \\
\midrule
Umbral HU & $>-200$ & Separa tejido de aire \\
Elemento estructurante & $3\times3\times3$ vox. & Cubo para cierre morfologico \\
Iteraciones cierre & 3 & Relleno de huecos internos \\
Conectividad & 8-vecina & Vecindario 2D por slice \\
Fondo & $-1024$ HU & Valor de aire asignado \\
\bottomrule
\end{tabular}
\end{table}

\subsection{Control de Calidad (AI Supervisor)}

\begin{table}[H]
\centering
\begin{tabular}{llc}
\toprule
\textbf{Verificacion} & \textbf{Metrica} & \textbf{Umbral} \\
\midrule
Volumen eliminado & $f_{\text{rem}} = (V_{\text{orig}} - V_{\text{masc}})/V_{\text{orig}}$ & $< 0.5$ (50\%) \\
Mascara no vacia & $\sum M \geq 1$ voxel & Critico \\
Nodo CT\_sin\_camilla & Existe en escena Slicer & Critico \\
\bottomrule
\end{tabular}
\end{table}

\newpage

% ============================================================
% 5. REGISTRO PET/CT
% ============================================================
\section{Registro y Re-muestreo PET/CT}
\label{sec:registro}

\subsection*{Acronimos usados en esta seccion}
\begin{tabular}{ll}
\toprule
\textbf{Acronimo} & \textbf{Significado} \\
\midrule
ASGD & Adaptive Stochastic Gradient Descent --- optimizador de Elastix \\
BSpline & Basis Spline --- funcion de interpolacion polinomica \\
CLI & Command-Line Interface \\
DOF & Degrees of Freedom --- grados de libertad \\
MI & Mutual Information --- metrica de similitud \\
MSE & Mean Squared Error \\
NCC & Normalized Cross-Correlation \\
RMSE & Root Mean Squared Error \\
\bottomrule
\end{tabular}

El PET y el CT se adquieren en equipos diferentes con geometrias distintas. Sin este paso, no
existe correspondencia voxel a voxel entre la actividad metabolica (PET) y la anatomia (CT).
El PET (tipicamente $200\times200\times N$ con espaciado $\sim 4.07$ mm) se re-muestrea a la
grilla del CT ($512\times512\times N$ con $\sim 0.97$ mm), conservando la actividad total.
El pipeline implementa tres metodos complementarios, con el primero como default.

\subsection{Problema Geometrico}

El PET DICOM y el CT tienen geometrias fundamentalmente distintas:

\begin{table}[H]
\centering
\begin{tabular}{l c c}
\toprule
\textbf{Parametro} & \textbf{PET} & \textbf{CT} \\
\midrule
Dimensiones ($N_x \times N_y$) & $200 \times 200$ & $512 \times 512$ \\
Espaciado ($s_x, s_y$) [mm] & $4.07 \times 4.07$ & $0.976 \times 0.976$ \\
Espaciado axial ($s_z$) [mm] & $2.0$ & $3.0$ \\
Voxels por slice & 40,000 & 262,144 \\
\bottomrule
\end{tabular}
\end{table}

\subsection{Notacion General}

\begin{table}[H]
\centering
\begin{tabular}{lll}
\toprule
\textbf{Variable} & \textbf{Descripcion} & \textbf{Unidad} \\
\midrule
$V_{\text{PET}}$ & Volumen PET en coordenadas voxel IJK & --- \\
$V_{\text{CT}}$ & Volumen CT en coordenadas voxel IJK & --- \\
$T_{\text{PET}}$ & Transformacion IJK$\to$RAS del PET ($4\times4$) & --- \\
$T_{\text{CT}}$ & Transformacion IJK$\to$RAS del CT ($4\times4$) & --- \\
$s_x^{\text{CT}}, s_y^{\text{CT}}, s_z^{\text{CT}}$ & Espaciado del CT & mm \\
$\mathbf{r}_{\text{CT}}$ & Coordenada RAS de un voxel CT & mm \\
$o_x^{\text{PET}}, o_y^{\text{PET}}, o_z^{\text{PET}}$ & Origen RAS del PET & mm \\
\bottomrule
\end{tabular}
\end{table}

\subsection{Metodo A --- BRAINSResample (Pipeline por Defecto)}

Usa el modulo CLI \texttt{BRAINSResample} de Slicer:

\begin{verbatim}
inputVolume     = PET_node
referenceVolume = CT_node
outputVolume    = PET_registrado
interpolationType = "Linear"
pixelType       = "float"
\end{verbatim}

\subsubsection{Opciones de Interpolacion}

\begin{table}[H]
\centering
\begin{tabular}{l l l}
\toprule
\textbf{Tipo} & \textbf{Descripcion} & \textbf{Uso tipico} \\
\midrule
\texttt{NearestNeighbor} & Asigna el valor del voxel mas cercano & Labelmaps \\
\texttt{Linear} & Interpolacion trilineal (8 vecinos) & PET, CT (default) \\
\texttt{BSpline} & Spline cubico de 3er orden (64 vecinos) & PET con suavizado \\
\texttt{Cubic} & Convolucion cubica (16 vecinos) & Aplicaciones cientificas \\
\bottomrule
\end{tabular}
\end{table}

\subsubsection{Algoritmo}

\begin{enumerate}
    \item Para cada voxel $\mathbf{r}_{\text{CT}}$ en la grilla destino (CT), calcular su
          equivalente en coordenadas PET:
          \begin{equation}\label{eq:coord_pet}
              \mathbf{r}_{\text{PET}} = T_{\text{PET}}^{-1} \cdot T_{\text{CT}} \cdot \mathbf{r}_{\text{CT}}
          \end{equation}
    \item Si $\mathbf{r}_{\text{PET}}$ cae fuera del volumen PET, asignar 0.
    \item Si cae dentro, interpolar segun el metodo elegido.
    \item Clipear negativos a 0 y aplicar conservacion de actividad:
          \begin{equation}\label{eq:act_post}
              \text{PET}_{\text{final}} = \text{PET}_{\text{interp}} \cdot \frac{A_{\text{pre-resample}}}{A_{\text{post}}}
          \end{equation}
\end{enumerate}

\subsection{Metodo B --- NumPy map\_coordinates}

Alternativa mas exacta que replica el algoritmo MATLAB \texttt{register\_v7.m}:

\begin{enumerate}
    \item Extraer arrays numpy de PET y CT.
    \item Calcular grillas mundo para cada imagen:
          \begin{equation}\label{eq:grilla_mundo}
              x_{\text{PET}}(i) = i \cdot s_x^{\text{PET}} + o_x^{\text{PET}}
          \end{equation}
    \item Interpolar con \texttt{scipy.ndimage.map\_coordinates} con \texttt{order=1}.
    \item Aplicar conservacion de actividad.
\end{enumerate}

\subsection{Metodo C --- Elastix (Registro Espacial)}

El modulo \texttt{DosimetryRegistration} en \texttt{registration.py} provee registro espacial
elastico (no solo remuestreo):

\begin{table}[H]
\centering
\begin{tabular}{l l l l}
\toprule
\textbf{Metodo} & \textbf{Transformacion} & \textbf{DOF} & \textbf{Optimizador} \\
\midrule
BrainsFit & Rigid + Affine + BSpline & $6 + 6 + \text{vars}$ & 1500 iteraciones \\
Elastix rigido & Euler (3T + 3R) & 6 & MI + ASGD \\
Elastix affine & Rigido + escala + shear & 12 & MI + ASGD \\
Elastix BSpline & Rigido + BSpline no rigido & variables & MI + ASGD \\
\bottomrule
\end{tabular}
\end{table}

donde:
\begin{itemize}
    \item DOF: Degrees of Freedom (grados de libertad).
    \item 3T: 3 traslaciones ($T_x, T_y, T_z$).
    \item 3R: 3 rotaciones ($R_x, R_y, R_z$).
\end{itemize}

\subsection{Metricas de Similitud}

\begin{table}[H]
\centering
\begin{tabular}{l l l}
\toprule
\textbf{Metrica} & \textbf{Formula} & \textbf{Rango} \\
\midrule
NCC & $\displaystyle\frac{\sum (A-\bar{A})(B-\bar{B})}
                {\sqrt{\sum (A-\bar{A})^2 \sum (B-\bar{B})^2}}$ & $[-1, 1]$ \\
MSE & $\displaystyle\frac{1}{N}\sum (A-B)^2$ & $[0, \infty)$ \\
MI & $\displaystyle\sum_{a,b} p(a,b)\log\frac{p(a,b)}{p(a)p(b)}$ & $[0, \infty)$ \\
MAE & $\displaystyle\frac{1}{N}\sum |A-B|$ & $[0, \infty)$ \\
RMSE & $\sqrt{\text{MSE}}$ & $[0, \infty)$ \\
\bottomrule
\end{tabular}
\end{table}

\subsection{Control de Calidad (AI Supervisor)}

\begin{table}[H]
\centering
\begin{tabular}{llc}
\toprule
\textbf{Verificacion} & \textbf{Metrica} & \textbf{Umbral} \\
\midrule
Alineacion PET/CT & NCC entre PET y CT registrados & $> 0.6$ \\
Conservacion actividad & $|A_{\text{final}}/A_{\text{original}} - 1|$ & $< 1\%$ \\
Dimensiones post-registro & $= 512 \times 512 \times N_{\text{CT}}$ & Critico \\
Negativos post-interpolacion & $N_{\text{voxels negativos}}$ & $= 0$ \\
\bottomrule
\end{tabular}
\end{table}

\newpage

% ============================================================
% 6. FUSION CT+PET
% ============================================================
\section{Fusion CT+PET y Dialogo de Actividad}
\label{sec:fusion}

\subsection*{Acronimos usados en esta seccion}
\begin{tabular}{ll}
\toprule
\textbf{Acronimo} & \textbf{Significado} \\
\midrule
CT & Computed Tomography \\
GBq & Gigabecquerel ($10^9$ Bq) \\
HU & Hounsfield Unit \\
mCi & Milicurie ($3.7 \times 10^7$ Bq) \\
NCC & Normalized Cross-Correlation \\
PET & Positron Emission Tomography \\
\bottomrule
\end{tabular}

Despues del registro, el PET esta en la misma geometria que el CT. La fusion visual permite al
medico verificar la alineacion y examinar la distribucion de actividad, mientras que el dialogo
informativo proporciona los datos cuantitativos de actividad total que se usaran en el calculo
dosimetrico. Este paso configura las vistas medicas de Slicer en layout 4-up y muestra un dialogo
modal con informacion del paciente, metadatos de imagen y actividad calculada desde DICOM raw.

\subsection{Layout de Visualizacion 4-up}

\begin{verbatim}
+-----------------------+-----------------------+
|        AXIAL          |       SAGITAL         |
|   (corte transversal) |   (corte sagital)     |
|   Background: CT      |   Background: CT      |
|   Foreground: PET     |   Foreground: PET     |
+-----------------------+-----------------------+
|       CORONAL         |        3D             |
|   (corte coronal)     |   (volumen rendering) |
|   Background: CT      |   CT + PET +          |
|   Foreground: PET     |   segmentos)          |
+-----------------------+-----------------------+
\end{verbatim}

\subsection{Parametros de Visualizacion}

\begin{table}[H]
\centering
\begin{tabular}{l l}
\toprule
\textbf{Parametro} & \textbf{Valor} \\
\midrule
Layout & \texttt{SlicerLayoutFourUpView} \\
Background (fondo) & CT (o CT\_sin\_camilla) \\
Foreground (primer plano) & PET en Bq/mL \\
Opacidad foreground & 0.35 (35\%) \\
Window/Level CT & Level 40 HU, Width 400 HU (abdomen) \\
Window/Level PET & P2--P98 de valores $>0$ (auto) \\
Colormap PET & Rainbow invertido (azul=bajo, rojo=alto) \\
Volume Rendering 3D & RayCast para CT y PET \\
Slice linking & Activado (navegacion sincronizada) \\
\bottomrule
\end{tabular}
\end{table}

\subsection{Window/Level}

La transformacion de valores HU a escala de grises:

\begin{equation}\label{eq:window_level}
    I_{\text{display}}(x) = \begin{cases}
        0 & \text{si } x < L - W/2 \\
        255 \cdot \dfrac{x - L + W/2}{W} & \text{si } L - W/2 \leq x \leq L + W/2 \\
        255 & \text{si } x > L + W/2
    \end{cases}
\end{equation}
\vardef{$I_{\text{display}}(x)$}{Valor de pixel en pantalla (escala de grises)}
\vardef{$L$}{Level (centro de la ventana). Para abdomen: $L=40$ HU}
\vardef{$W$}{Width (ancho de la ventana). Para abdomen: $W=400$ HU}

\subsection{Dialogo de Fusion (Fusion Dialog)}

Dialogo \textbf{modal} con las siguientes secciones informativas:

\begin{enumerate}
    \item \textbf{Paciente}: PatientID, peso (desde \texttt{pipeline\_config.jsonc}).
    \item \textbf{CT}: dimensiones, espaciado, nombre del nodo Slicer.
    \item \textbf{PET (DICOM raw)}: RescaleType, RescaleSlope, RescaleIntercept,
          numero de slices, volumen del voxel.
    \item \textbf{Actividad}: Total Bq, GBq, mCi, Bq/mL media y maxima, voxeles activos.
    \item \textbf{Registro PET$\to$CT}: metodo usado, duracion, verificacion de conservacion.
\end{enumerate}

\subsection{Verificaciones de Calidad}

\begin{table}[H]
\centering
\begin{tabular}{llc}
\toprule
\textbf{Verificacion} & \textbf{Metrica} & \textbf{Umbral} \\
\midrule
Unidades PET & \texttt{RescaleType == ``BQML''} & Critico \\
Rango de actividad & $0.1 \leq A_{\text{total}} \leq 50$ GBq & Warning \\
Dimensiones CT/PET & Similares tras registro & Warning \\
Voxeles activos & $N_{\text{activos}} > 100$ & Warning \\
Conservacion actividad & $|A_{\text{post}}/A_{\text{pre}} - 1| < 0.01$ & Warning \\
\bottomrule
\end{tabular}
\end{table}

\subsection{Control de Calidad (AI Supervisor)}

\begin{table}[H]
\centering
\begin{tabular}{llc}
\toprule
\textbf{Verificacion} & \textbf{Metrica} & \textbf{Umbral} \\
\midrule
Dialogo mostrado & Interaccion del usuario & Critico \\
Datos de actividad & Valores consistentes con DICOM raw & Warning \\
\bottomrule
\end{tabular}
\end{table}

\newpage

% ============================================================
% 7. ANONIMIZACION
% ============================================================
\section{Anonimizacion DICOM}
\label{sec:anonimizacion}

\subsection*{Acronimos usados en esta seccion}
\begin{tabular}{ll}
\toprule
\textbf{Acronimo} & \textbf{Significado} \\
\midrule
DICOM & Digital Imaging and Communications in Medicine \\
GDPR & General Data Protection Regulation (Union Europea) \\
HIPAA & Health Insurance Portability and Accountability Act (EEUU) \\
PHI & Protected Health Information --- datos de salud protegidos \\
TS & TotalSegmentator \\
\bottomrule
\end{tabular}

Las imagenes DICOM contienen informacion de salud protegida (PHI) que debe eliminarse antes de
procesarlas con herramientas externas como TotalSegmentator. Este paso copia los archivos DICOM
originales a un directorio \texttt{anon/}, limpia los metadatos identificables (nombre del
paciente, ID, fecha del estudio, institucion, etc.) y renombra los archivos con prefijo
\texttt{ANON\_}. Es un requisito legal segun HIPAA (EEUU) y GDPR (Europa).

\subsection{Campos Eliminados}

\begin{table}[H]
\centering
\begin{tabular}{l l l}
\toprule
\textbf{Tag DICOM} & \textbf{Campo} & \textbf{Razon} \\
\midrule
(0010,0010) & PatientName & Identificacion directa \\
(0010,0020) & PatientID & Identificacion directa \\
(0008,0020) & StudyDate & Informacion temporal \\
(0008,0080) & InstitutionName & Origen del estudio \\
(0008,0090) & ReferringPhysician & Medico tratante \\
(0008,0050) & AccessionNumber & ID interno del estudio \\
(0020,0010) & StudyID & ID del estudio \\
(0008,103E) & SeriesDescription & Descripcion de la serie \\
(0008,0081) & InstitutionAddress & Direccion de la institucion \\
\bottomrule
\end{tabular}
\end{table}

\subsection{Algoritmo}

\begin{enumerate}
    \item Crear directorio \texttt{anon/} dentro de \texttt{output\_dir}.
    \item Copiar los archivos DICOM de CT y PET al directorio anonimizado.
    \item Para cada archivo:
          \begin{itemize}
              \item Abrir con \texttt{pydicom.dcmread()}.
              \item Para cada tag en la lista de campos a limpiar, asignar string vacio.
              \item Guardar archivo modificado.
          \end{itemize}
    \item Renombrar archivos con prefijo \texttt{ANON\_}.
\end{enumerate}

\subsection{Control de Calidad (AI Supervisor)}

\begin{table}[H]
\centering
\begin{tabular}{llc}
\toprule
\textbf{Verificacion} & \textbf{Metrica} & \textbf{Umbral} \\
\midrule
Archivos anonimizados & \texttt{os.path.exists(anon\_ct, anon\_pet)} & Critico \\
Campos PHI limpiados & PatientName vacio & Warning \\
\bottomrule
\end{tabular}
\end{table}

\newpage

% ============================================================
% 8. TOTALSEGMENTATOR
% ============================================================
\section{TotalSegmentator --- Segmentacion Anatomica}
\label{sec:ts}

\subsection*{Acronimos usados en esta seccion}
\begin{tabular}{ll}
\toprule
\textbf{Acronimo} & \textbf{Significado} \\
\midrule
CT & Computed Tomography \\
GPU & Graphics Processing Unit \\
IA & Inteligencia Artificial \\
nnU-Net & Red convolucional para segmentacion 3D \\
TS & TotalSegmentator \\
\bottomrule
\end{tabular}

Segmentar manualmente 104 estructuras anatomicas en un CT abdominal tomaria horas a un radiologo.
TotalSegmentator (TS) automatiza completamente esta tarea usando una red neuronal nnU-Net
entrenada en miles de tomografias. Este paso ejecuta TS sobre el CT anonimizado, produciendo
segmentaciones individuales para cada organo, hueso, musculo y vaso, que luego seran mapeadas
a materiales MCNP.

\subsection{Ejecucion}

Se usa la API directa \texttt{TotalSegmentatorLogic.process()}:

\begin{verbatim}
from TotalSegmentator import TotalSegmentatorLogic
logic = TotalSegmentatorLogic()
logic.setupPythonRequirements()
logic.process(
    inputVolume=ct_node,
    outputSegmentation=seg_node,
    fast=True, cpu=True, task="total",
)
\end{verbatim}

\textbf{Importante}: TS es un ScriptedLoadableModule, NO un CLI module.
NO se usa \texttt{slicer.cli.run()}.

\subsection{Tasks Disponibles}

\begin{table}[H]
\centering
\begin{tabular}{l l c c}
\toprule
\textbf{Task} & \textbf{Descripcion} & \textbf{Clases} & \textbf{Soporta --fast} \\
\midrule
\texttt{total} & Todos los organos y estructuras & 104 & Si \\
\texttt{body} & Contorno externo del cuerpo & 1 & Si \\
\texttt{body\_lowdose} & Body para CT de baja dosis & 1 & Si \\
\texttt{liver\_lesions} & Lesiones hepaticas (tumores) & 1 & No \\
\texttt{lung\_vessels} & Vasos pulmonares & 1 & No \\
\bottomrule
\end{tabular}
\end{table}

\subsection{Clases Anatomicas Producidas}

\begin{description}
    \item[Organos]: higado, pulmones (5 lobulos), corazon, rinones, bazo, pancreas,
          vesicula, estomago, duodeno, intestino delgado, colon, vejiga, prostata,
          tiroides, adrenales, traquea, esofago.
    \item[Huesos]: vertebras (C1-C7, T1-T12, L1-L5, S1), costillas, esternon,
          claviculas, escapulas, humeros, femures, pelvis, craneo.
    \item[Musculos]: gluteos, iliopsoas, autochthon.
    \item[Vasos]: aorta, vena cava, arterias/venas iliacas, vena porta.
    \item[Cerebro]: cerebro, cerebelo, tronco encefalico.
    \item[Otros]: medula espinal, tejido subcutaneo, glandulas salivales, cristalinos.
\end{description}

\subsection{Tiempos de Ejecucion}

\begin{table}[H]
\centering
\begin{tabular}{l c c c}
\toprule
\textbf{Modo} & \textbf{GPU} & \textbf{CPU} & \textbf{Tiempo} \\
\midrule
fast & Si & --- & 2--5 min \\
fast & --- & Si & 5--15 min \\
full & Si & --- & 10--20 min \\
full & --- & Si & 30--60 min \\
\bottomrule
\end{tabular}
\end{table}

\subsection{Compatibilidad Python 3.9}

Slicer 5.8.1 usa Python 3.9.10. TotalSegmentator v2.13.0 requiere tres parches:

\begin{table}[H]
\centering
\begin{tabular}{c l l}
\toprule
\textbf{\#} & \textbf{Archivo} & \textbf{Parche} \\
\midrule
1 & \texttt{totalsegmentator/dicom\_io.py} & Agregar \texttt{from \_\_future\_\_ import annotations} \\
2 & \texttt{nnunetv2/.../data\_loader.py} & Shim \texttt{nnUNetDataLoader} $\to$ \texttt{nnUNetDataLoaderBase} \\
3 & \texttt{acvl\_utils} & Forzar version 0.2.5 \\
\bottomrule
\end{tabular}
\end{table}

\subsection{Control de Calidad (AI Supervisor)}

\begin{table}[H]
\centering
\begin{tabular}{llc}
\toprule
\textbf{Verificacion} & \textbf{Metrica} & \textbf{Umbral} \\
\midrule
Segmentos generados & $N_{\text{seg}} \geq 20$ & Critico \\
Higado presente & Segmento ``liver'' existe & Critico \\
Pulmones presentes & Segmentos pulmonares existen & Warning \\
Tiempo de ejecucion & $< 60$ min & Warning \\
\bottomrule
\end{tabular}
\end{table}

\newpage

% ============================================================
% 9. MAPEO TS -> PHANTOM
% ============================================================
\section{Mapeo TS \texorpdfstring{$\to$}{a} Phantom 3Dosim}
\label{sec:mapeo}

\subsection*{Acronimos usados en esta seccion}
\begin{tabular}{ll}
\toprule
\textbf{Acronimo} & \textbf{Significado} \\
\midrule
ICRP-110 & Publicacion 110: phantoms computacionales de referencia \\
ICRU-44 & Reporte 44: composiciones de tejidos para dosimetria \\
MCNP & Monte Carlo N-Particle \\
TS & TotalSegmentator \\
\bottomrule
\end{tabular}

TotalSegmentator produce 104 clases anatomicas con etiquetas numericas (labels 1--126). El
phantom 3Dosim usa un sistema de indices diferente ($30$=tejido blando, $50$=pulmon, $80$=hueso,
$90$=higado, $100$=tumor, $101$--$153$=materiales ICRP-110). Este paso convierte la segmentacion
TS a la representacion interna del pipeline mediante una tabla de lookup
(\texttt{ts\_label\_to\_phantom} en \texttt{tissue\_config.json}), asignando a cada voxel el
indice correcto para la generacion de entrada MCNP (Modulo 2).

\subsection{Tabla de Indices del Phantom}

\begin{table}[H]
\centering
\begin{tabular}{c l c c}
\toprule
\textbf{Indice} & \textbf{Tejido} & \textbf{Densidad [g/cm$^3$]} & \textbf{Material MCNP} \\
\midrule
1 & Aire & 0.001205 & M1 \\
30 & Tejido blando & 1.06 & M2 \\
50 & Pulmon & 0.382 & M3 \\
80 & Hueso & 1.92 & M4 \\
90 & Higado & 1.05 & M5 \\
100 & Tumor & --- & M6 \\
101--153 & ICRP-110 (53 materiales) & varios & M7--M59 \\
200 & Peritumoral & --- & M60 \\
\bottomrule
\end{tabular}
\end{table}

\subsection{Estructura de tissue\_config.json}

El archivo \texttt{tissue\_config.json} contiene tres secciones principales:

\subsubsection{ts\_label\_to\_phantom}

Diccionario que mapea labels numericos de TS (1--126) a indices del phantom:

\begin{verbatim}
"ts_label_to_phantom": {
    "1":  {"index": 30, "segment": "spleen"},
    "2":  {"index": 50, "segment": "lung_upper_lobe_left"},
    "5":  {"index": 90, "segment": "liver"},
    "8":  {"index": 80, "segment": "heart"},
    ...
}
\end{verbatim}

\subsubsection{ts\_body\_labels}

Lista de labels TS que corresponden a tejido corporal (no organos especificos).
Se asignan al indice 30 (Tejido blando).

\subsubsection{tissues}

Lista completa de 53+ tejidos con: indice, nombre, densidad, rango HU, color, composicion MCNP.

\subsection{Ejemplo de Composicion MCNP}

Para Higado (indice 90, densidad 1.05 g/cm$^3$):

\begin{verbatim}
"mcnp_material": {
    "id": 5,
    "composition": {
        "1000": 0.105,    // Hidrogeno (H)
        "6000": 0.143,    // Carbono (C)
        "7000": 0.034,    // Nitrogeno (N)
        "8000": 0.708,    // Oxigeno (O)
        "11000": 0.003,   // Sodio (Na)
        "15000": 0.003,   // Fosforo (P)
        "16000": 0.003,   // Azufre (S)
        "17000": 0.002,   // Cloro (Cl)
        "19000": 0.002    // Potasio (K)
    }
}
\end{verbatim}

\subsection{Algoritmo de Mapeo}

\begin{verbatim}
function _map_ts_to_phantom(ts_labelmap):
    1. Inicializar phantom = 1 (Aire) para todos los voxeles
    2. Asignar 30 (Tejido_blando) donde ts_labelmap esta en body_labels
    3. Para cada entrada en ts_label_to_phantom:
         mascara = (ts_labelmap == ts_label)
         phantom[mascara] = indice_phantom
    4. Retornar phantom
\end{verbatim}

\subsection{Correcciones de Bugs (Julio 2026)}

Se detectaron y corrigieron errores donde tejidos blandos estaban asignados a Hueso (densidad
1.92 g/cm$^3$) en vez de sus materiales ICRP-110 correctos:

\begin{table}[H]
\centering
\begin{tabular}{c l c c}
\toprule
\textbf{Label TS} & \textbf{Estructura} & \textbf{Antes (BUG)} & \textbf{Ahora (CORRECTO)} \\
\midrule
63--68 & Vasos sanguineos & Hueso (80) & Sangre (128) \\
79 & Medula espinal & Hueso (80) & Tejido mixto (145) \\
80--89 & Musculos & Hueso (80) & Tejido muscular (129) \\
90 & Cerebro & Hueso (80) & Cerebro (132) \\
15 & Esofago & Pulmon (50) & Esofago (144) \\
\bottomrule
\end{tabular}
\end{table}

Impacto: 22 estructuras anatomicas mal asignadas, causando sobre-estimacion de atenuacion y
distorsion de la dosis MCNP.

\subsection{Clase TissueConfig}

Singleton que carga y cachea \texttt{tissue\_config.json}:

\begin{table}[H]
\centering
\begin{tabular}{l l}
\toprule
\textbf{Metodo} & \textbf{Descripcion} \\
\midrule
get\_tissue(index) & Datos completos del tejido \\
get\_tissue\_density(index) & Densidad en g/cm$^3$ \\
get\_tissue\_hu\_range(index) & Rango de HU caracteristico \\
get\_mcnp\_material(index) & Composicion elemental MCNP \\
generate\_mcnp\_material\_card(index) & Genera tarjeta M para MCNP \\
get\_ts\_mapping() & Dict: label TS $\to$ indice phantom \\
get\_body\_labels() & Set de labels corporales \\
\bottomrule
\end{tabular}
\end{table}

\subsection{Control de Calidad (AI Supervisor)}

\begin{table}[H]
\centering
\begin{tabular}{llc}
\toprule
\textbf{Verificacion} & \textbf{Metrica} & \textbf{Umbral} \\
\midrule
Labels TS mapeados & Sin labels sin correspondencia & Warning \\
Densidad vs HU & Discrepancia $< 10\%$ & Warning \\
Integridad phantom & Todos los indices esperados presentes & Critico \\
\bottomrule
\end{tabular}
\end{table}

\newpage

% ============================================================
% 10. CREACION DE TUMOR
% ============================================================
\section{Creacion de Tumor (4 modos)}
\label{sec:tumor}

\subsection*{Acronimos usados en esta seccion}
\begin{tabular}{ll}
\toprule
\textbf{Acronimo} & \textbf{Significado} \\
\midrule
CT & Computed Tomography \\
IA & Inteligencia Artificial \\
MCNP & Monte Carlo N-Particle \\
NIfTI & Neuroimaging Informatics Technology Initiative \\
ROI & Region of Interest \\
TS & TotalSegmentator \\
\bottomrule
\end{tabular}

Para calcular la dosis en el tumor hepatico es necesario definir su contorno. El pipeline soporta
4 modos de creacion de tumor, desde una esfera sintetica hasta segmentacion automatica con IA.
Cada modo esta disenado para un escenario clinico diferente y se configura desde
\texttt{pipeline\_config.jsonc}. Opcionalmente se crea un segmento de ``higado sano'' como
diferencia entre el higado completo y el tumor.

\subsection{Configuracion}

\begin{verbatim}
"tumor": {
    "mode": "synthetic",                // synthetic | load_file |
                                        // manual | ts_liver_lesions
    "synthetic_radius_mm": 10.0,        // radio de la esfera (modo synthetic)
    "liver_segment_name": "liver",      // nombre del segmento hepatico en TS
    "create_healthy_liver": true,       // crear higado_sano = higado - tumor
    "load_file_path": "...",            // ruta NIfTI (modo load_file)
    "load_segment_name": "Tumor_Cargado",
    "manual_segment_name": "Tumor_Manual",
    "ts_liver_lesions_segment_name": "Tumor_TS",
    "ts_liver_lesions_min_volume_cc": 1.0,  // volumen minimo de lesion
    "timeout_minutes": 0                // timeout (0 = sin limite)
}
\end{verbatim}

\subsection{Modo 1: Synthetic (esfera sintetica)}

Crea una esfera de radio $R$ mm dentro del parenquima hepatico.

\subsubsection{Algoritmo}

\begin{enumerate}
    \item Extraer mascara del higado $M_{\text{higado}}$ desde TS.
    \item Calcular centroide del higado en coordenadas voxel:
          \begin{equation}\label{eq:centroide}
              c_i = \frac{1}{N} \sum_{n=1}^{N} i_n,\quad i \in \{x, y, z\}
          \end{equation}
          \vardef{$c_x, c_y, c_z$}{Coordenadas del centroide del higado}
          \vardef{$N$}{Numero total de voxeles hepaticos}
    \item Si el centroide cae fuera del higado, buscar el voxel hepatico mas cercano.
    \item Crear mascara esferica con distancia euclidiana en mm:
          \begin{equation}\label{eq:esfera}
              M_{\text{esfera}}(x,y,z) = \mathbb{1}\left[
              d_{\text{mm}}\big((x,y,z), (c_x,c_y,c_z)\big) \leq R
              \right]
          \end{equation}
          donde:
          \begin{equation}\label{eq:dist_mm}
              d_{\text{mm}}(\mathbf{r}_1, \mathbf{r}_2) = \sqrt{
              (x_1-x_2)^2 s_x^2 + (y_1-y_2)^2 s_y^2 + (z_1-z_2)^2 s_z^2}
          \end{equation}
          \vardef{$d_{\text{mm}}$}{Distancia euclidiana en mm}
          \vardef{$s_x, s_y, s_z$}{Espaciado del CT en mm}
          \vardef{$R$}{Radio de la esfera (configurable)}
    \item Intersectar con el higado:
          \begin{equation}\label{eq:tumor_mask}
              M_{\text{tumor}} = M_{\text{esfera}} \cap M_{\text{higado}}
          \end{equation}
    \item Agregar como segmento ``Tumor\_Sintetico'' (color rojo).
\end{enumerate}

\subsection{Modo 2: Load File (cargar NIfTI)}

\begin{enumerate}
    \item Cargar NIfTI con \texttt{slicer.util.loadNodeFromFile()}.
    \item Convertir a mascara binaria.
    \item Si las dimensiones no coinciden con el CT, re-muestrear con BRAINSResample
          (interpolacion NearestNeighbor).
    \item Enmascarar dentro del higado: $M_{\text{tumor}} = M_{\text{NIfTI}} \cap M_{\text{higado}}$.
    \item Agregar como segmento con nombre configurable.
\end{enumerate}

\subsection{Modo 3: Manual (Segment Editor)}

\begin{enumerate}
    \item Pipeline crea segmento vacio ``Tumor\_Manual''.
    \item Activa modulo \texttt{Segment Editor} de Slicer.
    \item Muestra dialogo NO MODAL con instrucciones.
    \item Medico dibuja el tumor (Paint, Scissors, Level Tracing, GrowCut).
    \item Medico hace clic en APROBAR.
\end{enumerate}

\subsection{Modo 4: TS Liver Lesions (automatico con IA)}

\begin{enumerate}
    \item Ejecutar TS con \texttt{task=``liver\_lesions''} sobre el CT.
          \textbf{Nota}: este task NO soporta modo \texttt{--fast}.
    \item Cargar resultado como mascara binaria.
    \item Enmascarar dentro del higado (phantom idx 90).
    \item Post-procesamiento: componentes conectadas, eliminar volumen $< V_{\min}$.
    \item Agregar como segmento ``Tumor\_TS''.
\end{enumerate}

\subsection{Higado Sano}

En todos los modos, opcionalmente se crea ``higado\_sano'' como diferencia:

\begin{equation}\label{eq:higado_sano}
    M_{\text{higado\_sano}} = M_{\text{higado}} \setminus M_{\text{tumor}}
\end{equation}

\subsection{Control de Calidad (AI Supervisor)}

\begin{table}[H]
\centering
\begin{tabular}{llc}
\toprule
\textbf{Verificacion} & \textbf{Metrica} & \textbf{Umbral} \\
\midrule
Volumen del tumor & $V_{\text{tumor}} \geq 0.5$ cm$^3$ & Warning \\
Tumor dentro del higado & $M_{\text{tumor}} \subseteq M_{\text{higado}}$ & Critico \\
Tumor no vacio & $N_{\text{vox,tumor}} > 0$ & Critico \\
Lesiones detectadas (TS) & $N_{\text{lesiones}} \geq 1$ & Warning \\
\bottomrule
\end{tabular}
\end{table}

\newpage

% ============================================================
% 11. VALIDACION MEDICA
% ============================================================
\section{Validacion Medica}
\label{sec:validacion}

\subsection*{Acronimos usados en esta seccion}
\begin{tabular}{ll}
\toprule
\textbf{Acronimo} & \textbf{Significado} \\
\midrule
CT & Computed Tomography \\
PET & Positron Emission Tomography \\
TS & TotalSegmentator \\
\bottomrule
\end{tabular}

Hay decisiones que ningun algoritmo puede tomar. La validacion medica es un punto de control
obligatorio donde un medico especialista debe aprobar explicitamente cada etapa critica del
pipeline antes de continuar. El pipeline tiene dos puntos de validacion: uno tras la
segmentacion anatomica (TotalSegmentator) y otro tras la creacion del tumor. Ambos usan
dialogos NO modales que permiten al medico navegar las imagenes libremente mientras decide.
No existe ``auto-approve'' bajo ninguna circunstancia.

\subsection{Validacion de Segmentacion}

Se ejecuta inmediatamente despues de TotalSegmentator:

\begin{verbatim}
Tras TotalSegmentator
    |
    v
setup_medical_views(segnode)  <- segmentaciones visibles en 3D
    |
    v
show_validation_dialog("segmentacion")
    |
    +-- APROBAR -> continua pipeline
    +-- RECHAZAR -> pipeline detenido (RuntimeError)
\end{verbatim}

\subsection{Validacion de Tumor}
\label{sec:validacion_tumor}

Se ejecuta despues de crear el tumor:

\begin{verbatim}
Tras crear tumor
    |
    v
setup_medical_views(segnode)
    |
    v
show_tumor_validation_dialog(context=mode)
    |
    +-- APROBAR -> continua (higado sano, body, labelmap)
    +-- RECHAZAR -> pipeline detenido
\end{verbatim}

\subsection{Caracteristicas del Dialogo}

\begin{table}[H]
\centering
\begin{tabular}{l l}
\toprule
\textbf{Caracteristica} & \textbf{Valor} \\
\midrule
Modalidad & NO modal (\texttt{setModal(False)}) \\
Loop de eventos & \texttt{slicer.app.processEvents()} \\
Overlay PET & Slider de opacidad (0--100\%) \\
Layout & 4 vistas (axial + sagital + coronal + 3D) \\
Texto instructivo & Contextual segun el paso \\
Botones & APROBAR (verde), RECHAZAR (rojo) \\
Consecuencia rechazo & \texttt{raise RuntimeError} \\
\bottomrule
\end{tabular}
\end{table}

\subsection{Control de Calidad (AI Supervisor)}

\begin{table}[H]
\centering
\begin{tabular}{llc}
\toprule
\textbf{Verificacion} & \textbf{Metrica} & \textbf{Umbral} \\
\midrule
Respuesta del dialogo & APROBAR o RECHAZAR & Critico \\
Tiempo en dialogo & $< 30$ min & Warning \\
Validacion en checkpoint & Forzada incluso al restaurar & Critico \\
\bottomrule
\end{tabular}
\end{table}

\newpage

% ============================================================
% 12. BODY SEGMENTATION
% ============================================================
\section{Segmentacion Corporal (Body)}
\label{sec:body}

\subsection*{Acronimos usados en esta seccion}
\begin{tabular}{ll}
\toprule
\textbf{Acronimo} & \textbf{Significado} \\
\midrule
CT & Computed Tomography \\
MCNP & Monte Carlo N-Particle \\
TS & TotalSegmentator \\
\bottomrule
\end{tabular}

La labelmap dosimetrica necesita un contorno externo que defina donde termina el cuerpo del
paciente y comienza el aire. Este paso ejecuta TotalSegmentator con \texttt{task=``body''} para
obtener una mascara binaria del contorno externo, que luego se usa como ``contenedor'' en la
exportacion de labelmap: todos los voxeles dentro del cuerpo que no pertenecen a ningun organo
especifico se asignan como tejido blando (indice 30), garantizando que no queden voxeles sin
material asignado para la simulacion MCNP.

\subsection{Algoritmo}

\begin{enumerate}
    \item Ejecutar TotalSegmentator con \texttt{task=``body''} sobre el CT.
    \item Cargar resultado como nodo \texttt{Body\_Segmentation}.
    \item Si falla, continuar sin body (warning).
\end{enumerate}

\subsection{Integracion con la Labelmap}

En \texttt{labelmap\_exporter.py}, el body se incorpora al final:

\begin{enumerate}
    \item Primero se asignan todos los organos (higado, pulmon, etc.).
    \item Luego se asigna body como tejido blando (indice 30) donde no hay organos:
          \begin{equation}\label{eq:body_labelmap}
              \text{body\_region} = (M_{\text{body}} > 0) \land
              \lnot \text{any\_organ}
          \end{equation}
\end{enumerate}

Esto garantiza que no haya solapamiento: los organos tienen prioridad.

\subsection{Control de Calidad (AI Supervisor)}

\begin{table}[H]
\centering
\begin{tabular}{llc}
\toprule
\textbf{Verificacion} & \textbf{Metrica} & \textbf{Umbral} \\
\midrule
Body contiene organos & Ningun organo fuera de body\_mask & Critico \\
Voxels sin asignar & $N_{\text{sin-asignar}} = 0$ & Warning \\
Body no vacio & $\sum(\text{body\_mask}) > 0$ & Critico \\
\bottomrule
\end{tabular}
\end{table}

\newpage

% ============================================================
% 13. EXPORTAR LABELMAP
% ============================================================
\section{Exportacion de Labelmap Dosimetrica}
\label{sec:labelmap}

\subsection*{Acronimos usados en esta seccion}
\begin{tabular}{ll}
\toprule
\textbf{Acronimo} & \textbf{Significado} \\
\midrule
MCNP & Monte Carlo N-Particle \\
NIfTI & Neuroimaging Informatics Technology Initiative \\
NRRD & Nearly Raw Raster Data \\
TS & TotalSegmentator \\
\bottomrule
\end{tabular}

Convertir la segmentacion completa (organos + tumor + body) en un archivo NIfTI y NRRD donde
cada voxel tiene un valor entero que corresponde al indice del tejido en
\texttt{tissue\_config.json}. Esta labelmap es la entrada del Modulo 2 (generacion de entrada
MCNP).

\subsection{Algoritmo Detallado}

\subsubsection{Paso 1: Cargar tissue\_config.json}

Se busca el archivo en \texttt{SlicerDosim/Resources/Config/tissue\_config.json}
o path absoluto configurable.

\subsubsection{Paso 2: Construir mapping nombre \texorpdfstring{$\to$}{a} indice}

Para cada segmento $s$ en la segmentacion:

\begin{enumerate}
    \item Si el ID del segmento es un numero (label TS original), lookup en
          \texttt{ts\_label\_to\_phantom}.
    \item Si no, lookup por nombre en \texttt{DEFAULT\_NAME\_TO\_PHANTOM}.
    \item Si no, detectar hueso por palabras clave (\texttt{``rib''}, \texttt{``vertebra''},
          \texttt{``scapula''}, \texttt{``sternum''}, \texttt{``bone''}, \texttt{``costal''}).
    \item Si no, asignar indice libre 2--256.
\end{enumerate}

\subsubsection{Paso 3: Acumular en labelmap 3D}

\begin{verbatim}
final_labelmap = zeros((Nz, Ny, Nx), dtype=int16)
any_organ = zeros((Nz, Ny, Nx), dtype=bool)

for cada segmento s con indice p:
    mask = extraer_mascara(s)
    # Resolver solapamientos: gana el indice mas alto
    new_region = (mask > 0) & ((final_labelmap == 0) | (p > final_labelmap))
    final_labelmap[new_region] = p
    any_organ[new_region] = True
\end{verbatim}

\subsubsection{Paso 4: Incorporar body segmentation}

\begin{verbatim}
body_mask = extraer_mascara(body_segmentation_node)
body_region = (body_mask > 0) & not any_organ
final_labelmap[body_region] = 30  # Tejido_blando
\end{verbatim}

\subsubsection{Paso 5: Verificar integridad}

\begin{table}[H]
\centering
\begin{tabular}{l l}
\toprule
\textbf{Verificacion} & \textbf{Que detecta} \\
\midrule
Valores unicos & Todos los indices esperados presentes \\
Overlap voxels & Solapamiento entre organos (ideal: 0) \\
Voxels sin asignar & Dentro del body sin tejido asignado \\
Body vs organos & Body no solapa con ningun organo \\
\bottomrule
\end{tabular}
\end{table}

\subsubsection{Paso 6: Exportar}

\begin{table}[H]
\centering
\begin{tabular}{l l l}
\toprule
\textbf{Formato} & \textbf{Extension} & \textbf{Uso} \\
\midrule
NIfTI & .nii & Modulo 2 (MCNP) \\
NRRD & .nrrd & 3D Slicer (formato nativo) \\
\bottomrule
\end{tabular}
\end{table}

\subsection{Control de Calidad (AI Supervisor)}

\begin{table}[H]
\centering
\begin{tabular}{llc}
\toprule
\textbf{Verificacion} & \textbf{Metrica} & \textbf{Umbral} \\
\midrule
Overlap voxels & $N_{\text{overlap}} = 0$ & Critico \\
Voxels sin asignar & $N_{\text{sin-asignar}} = 0$ & Critico \\
Indices esperados & Todos presentes en labelmap & Warning \\
\bottomrule
\end{tabular}
\end{table}

\newpage

% ============================================================
% 14. CHECKPOINT SYSTEM
% ============================================================
\section{Checkpoint System --- Persistencia y Reanudacion}
\label{sec:checkpoint}

\subsection*{Acronimos usados en esta seccion}
\begin{tabular}{ll}
\toprule
\textbf{Acronimo} & \textbf{Significado} \\
\midrule
JSON & JavaScript Object Notation \\
MRB & Medical Reality Bundle (escena 3D Slicer) \\
SIGINT & Signal Interrupt (Ctrl+C) \\
TS & TotalSegmentator \\
\bottomrule
\end{tabular}

En un pipeline de dosimetria que puede ejecutarse durante 15--60 minutos, las interrupciones son
inevitables: el usuario cierra Slicer, la computadora se reinicia, una segmentacion falla por
falta de memoria. El sistema de checkpoints resuelve este problema guardando el estado completo
del pipeline despues de \textbf{cada paso exitoso} en un archivo JSON estructurado. Al reiniciar,
el orquestador \texttt{PipelineMod1} detecta automaticamente el ultimo checkpoint y reanuda desde
el primer paso pendiente, restaurando nodos, parametros y visualizacion sin perdida de progreso.

\subsection{Estructura del Archivo de Checkpoint}

Se guarda en \texttt{output\_dir/checkpoints/pipeline\_checkpoint.json}:

\begin{verbatim}
{
    "pipeline_name": "PipelineMod1",
    "version": 4,
    "start_time": "2026-07-02T10:30:00",
    "last_update": "2026-07-02T10:45:00",
    "steps": [
        {"name": "Carga_DICOM",       "status": "completed",
                                        "time": "10:32:00"},
        {"name": "TotalSegmentator",   "status": "completed",
                                        "time": "10:50:00"},
        {"name": "Creacion_Tumor",     "status": "in_progress",
                                        "time": "10:55:00"},
        {"name": "Exportar_Labelmap",  "status": "pending",
                                        "time": ""}
    ],
    "metadata": {
        "patient_id": "ANON_001",
        "ct_file": "CT_anon.nrrd",
        "pet_file": "PET_anon.nii",
        "segnode_file": "segmentation_initial.seg.nrrd"
    },
    "data": {
        "load_dicom": {
            "ct_node_name": "CT",
            "pet_node_name": "PET"
        }
    }
}
\end{verbatim}

\subsection{Propiedades Completas del Checkpoint}

\begin{table}[H]
\centering
\begin{tabular}{l l l}
\toprule
\textbf{Propiedad} & \textbf{Tipo} & \textbf{Descripcion} \\
\midrule
pipeline\_name & string & Identificador del pipeline \\
version & int & Control de compatibilidad entre versiones \\
start\_time & datetime & Fecha/hora de inicio original \\
last\_update & datetime & Actualizado en cada paso \\
steps & array & Array con estados individuales de cada paso \\
metadata & dict & Archivos intermedios para restauracion \\
data & dict & Datos de restauracion por paso (nombres de nodos, paths) \\
\bottomrule
\end{tabular}
\end{table}

\subsection{Estados de los Pasos}

\begin{table}[H]
\centering
\begin{tabular}{l l l}
\toprule
\textbf{Estado} & \textbf{Significado} & \textbf{Accion al reanudar} \\
\midrule
\texttt{completed} & Ejecutado exitosamente & Se salta al reanudar \\
\texttt{in\_progress} & En ejecucion o fallo & Se re-ejecuta al reanudar \\
\texttt{pending} & No iniciado & Se ejecuta normalmente \\
\bottomrule
\end{tabular}
\end{table}

\subsection{Flujo de Restauracion}

\begin{enumerate}
    \item Al iniciar el pipeline, verificar si existe \texttt{pipeline\_checkpoint.json}.
    \item Si no existe $\to$ empezar desde el paso 1.
    \item Si existe $\to$ leer el array \texttt{steps}.
    \item Buscar el primer paso con estado \texttt{pending} o \texttt{in\_progress}.
    \item Reanudar desde ese paso, manteniendo los pasos \texttt{completed}.
    \item Restaurar nodos Slicer desde \texttt{data} guardados en checkpoint.
    \item Restaurar visualizacion medica con \texttt{setup\_medical\_views()}.
\end{enumerate}

\subsection{Restauracion de Estado}

Cuando un paso ya esta completado, \texttt{\_checkpoint\_step()} salta la ejecucion pero llama a
\texttt{\_restore\_step\_state()} para que el pipeline tenga disponibles los nodos necesarios:

\begin{table}[H]
\centering
\begin{tabular}{l l}
\toprule
\textbf{Data key} & \textbf{Attr a restaurar} \\
\midrule
ct\_node\_name & ct\_node (busca por nombre en escena) \\
pet\_node\_name & pet\_node \\
segmentation\_node\_name & segmentation\_node \\
ct\_masked\_node\_name & ct\_masked\_node \\
\bottomrule
\end{tabular}
\end{table}

\subsection{Guardado de Escena}

El pipeline guarda la escena \texttt{.mrb} independientemente del checkpoint JSON:

\begin{table}[H]
\centering
\begin{tabular}{l l}
\toprule
\textbf{Frecuencia} & \textbf{Puntos de guardado} \\
\midrule
\texttt{minimal} (default) & Post-carga DICOM y post-body segmentation \\
\texttt{all} & Tras cada paso \\
\bottomrule
\end{tabular}
\end{table}

\subsection{Control de Calidad (AI Supervisor)}

\begin{table}[H]
\centering
\begin{tabular}{llc}
\toprule
\textbf{Verificacion} & \textbf{Metrica} & \textbf{Umbral} \\
\midrule
Archivo JSON valido & \texttt{json.load()} exitoso & Critico \\
Version de checkpoint & \texttt{state[``version''] == CHECKPOINT\_VERSION} & Critico \\
Nodos restaurados & \texttt{slicer.util.getNode()} exitoso & Warning \\
Escena .mrb cargable & \texttt{slicer.util.loadScene()} exitoso & Warning \\
Validacion medica forzada & Forzada incluso si checkpoint indica completed & Critico \\
\bottomrule
\end{tabular}
\end{table}

\newpage

% ============================================================
% 15. AI SUPERVISOR
% ============================================================
\section{AI Supervisor --- Verificacion Automatica Post-Paso}
\label{sec:ai}

\subsection*{Acronimos usados en esta seccion}
\begin{tabular}{ll}
\toprule
\textbf{Acronimo} & \textbf{Significado} \\
\midrule
IA & Inteligencia Artificial \\
NCC & Normalized Cross-Correlation \\
TS & TotalSegmentator \\
\bottomrule
\end{tabular}

El AI Supervisor es un revisor inteligente que analiza el resultado de cada paso del pipeline
y detecta anomalias antes de continuar. No reemplaza la validacion medica, sino que la
complementa con verificaciones automaticas inmediatas. Utiliza reglas duras (pre-verificacion)
y opcionalmente consulta a un modelo de lenguaje (DeepSeek via OpenRouter) para analisis mas
profundo.

\subsection{Verificaciones por Paso --- Tabla Completa}

\begin{table}[H]
\centering
\begin{tabular}{l l l l}
\toprule
\textbf{Paso} & \textbf{Verificacion} & \textbf{Condicion de fallo} & \textbf{Accion} \\
\midrule
Carga DICOM & Dimensiones, modalidades & CT o PET no detectados & Error \\
Calibracion PET & Actividad total $A_{\text{total}}$ & $A_{\text{total}} \leq 0$ Bq & Error \\
Eliminar camilla & Fraccion eliminada $f_{\text{rem}}$ & $f_{\text{rem}} > 0.5$ & Warning \\
Registro PET$\to$CT & NCC entre PET y CT & NCC $< 0.6$ & Warning \\
TotalSegmentator & Segmentos $N_{\text{seg}}$ & $N_{\text{seg}} < 20$ & Error \\
Validacion medica & Respuesta del dialogo & Sin interaccion & Error \\
Creacion tumor & Volumen $V_{\text{tumor}}$ & $V_{\text{tumor}} < 0.5$ cm$^3$ & Warning \\
Body segmentation & Cobertura de organos & Organos fuera del body & Error \\
Export labelmap & Overlap voxels & Overlap $> 0$ & Error \\
\bottomrule
\end{tabular}
\end{table}

\subsection{Flujo de Operacion}

\begin{verbatim}
Fin de paso N
    |
    v
AI Supervisor: analizar resultado
    |
    +-- Normal -> continuar
    +-- Warning -> log + continuar
    +-- Error  -> detener pipeline
\end{verbatim}

\subsection{Configuracion}

\begin{verbatim}
"ai_supervisor": {
    "fail_on_warning": false,    // true = detener pipeline en warning
    "log_level": "warning",      // "info", "warning", "error"
    "checks": {
        "carga_dicom": true,
        "calibracion_pet": true,
        "couch_remover": true,
        "registro_pet_ct": true,
        "totalsegmentator": true,
        "creacion_tumor": true,
        "export_labelmap": true
    }
}
\end{verbatim}

\subsection{Pre-verificacion Inmediata (Reglas Duras)}

\begin{table}[H]
\centering
\begin{tabular}{l l l}
\toprule
\textbf{Regla} & \textbf{Que detecta} & \textbf{Accion} \\
\midrule
Metodo simple (threshold) & Segmentacion insuficiente & Warning \\
Pocos segmentos ($\leq 2$) & Threshold sin etiquetado de organos & Warning \\
Voxels fuera del cuerpo & Segmentacion incluye aire/camilla & Alerta \\
Actividad total cero & Calibracion PET fallida & Error \\
\bottomrule
\end{tabular}
\end{table}

\subsection{Control de Calidad (AI Supervisor)}

\begin{table}[H]
\centering
\begin{tabular}{llc}
\toprule
\textbf{Verificacion} & \textbf{Metrica} & \textbf{Umbral} \\
\midrule
Log de supervisor & \texttt{ai\_supervisor.log} escrito & Critico \\
Tiempo de respuesta IA & $< 30$ s (timeout DeepSeek) & Warning \\
Pre-verificaciones & Todas las reglas evaluadas & Critico \\
\bottomrule
\end{tabular}
\end{table}

\newpage

% ============================================================
% 16. VISUALIZACION 3D
% ============================================================
\section{Visualizacion 3D Automatica}
\label{sec:vis}

\subsection*{Acronimos usados en esta seccion}
\begin{tabular}{ll}
\toprule
\textbf{Acronimo} & \textbf{Significado} \\
\midrule
CT & Computed Tomography \\
HU & Hounsfield Unit \\
PET & Positron Emission Tomography \\
TS & TotalSegmentator \\
\bottomrule
\end{tabular}

Configurar la interfaz de 3D Slicer para mostrar las imagenes y segmentaciones en un layout
4-up optimizado para visualizacion medica: axial, sagital, coronal y vista 3D. La funcion
\texttt{setup\_medical\_views()} se llama tras cada paso critico del pipeline para que el
medico pueda navegar e inspeccionar los resultados sin configuracion manual.

\subsection{Momentos de Activacion}

\begin{table}[H]
\centering
\begin{tabular}{l l}
\toprule
\textbf{Momento} & \textbf{Proposito} \\
\midrule
Tras carga DICOM & Mostrar CT y PET iniciales \\
Tras registro PET$\to$CT & Mostrar PET re-muestreado sobre CT \\
Tras fusion CT+PET & Mostrar las 3 imagenes (CT, PET, fusion) \\
Tras TotalSegmentator & Mostrar segmentaciones sobre CT \\
Tras validacion & Mostrar resultado final \\
\bottomrule
\end{tabular}
\end{table}

\subsection{Funcion setup\_medical\_views()}

\begin{enumerate}
    \item Configurar layout 4-up:
          \texttt{layoutManager.setLayout(SlicerLayoutFourUpView)}.
    \item Sincronizar vistas (crosshair, posicion de slice).
    \item Configurar foreground y background en cada slice:
          \begin{itemize}
              \item Background: CT (window/level abdomen: $L=40$ HU, $W=400$ HU).
              \item Foreground: PET (opacidad 0.35, colormap rainbow invertido).
              \item Window/level PET: percentiles $P_2$ a $P_{98}$ de valores $>0$.
          \end{itemize}
    \item Activar volumen rendering 3D con RayCast para CT y PET.
    \item Resetear focal point al centro del paciente.
    \item Activar slice linking (navegacion sincronizada).
\end{enumerate}

\subsection{Control de Calidad (AI Supervisor)}

\begin{table}[H]
\centering
\begin{tabular}{llc}
\toprule
\textbf{Verificacion} & \textbf{Metrica} & \textbf{Umbral} \\
\midrule
Layout activo & Layout 4-up configurado & Critico \\
Slice linking & Navegacion sincronizada activada & Warning \\
Volume rendering & RayCast activo para CT y PET & Warning \\
\bottomrule
\end{tabular}
\end{table}

\newpage

% ============================================================
% ANEXO FUNDAMENTOS FISICOS
% ============================================================
\appendix
\section{Fundamentos Fisicos del \texorpdfstring{$^{90}$Y}{Y-90}}

\subsection*{Acronimos usados en este apendice}
\begin{tabular}{ll}
\toprule
\textbf{Acronimo} & \textbf{Significado} \\
\midrule
BED & Biologically Effective Dose \\
DVH & Dose-Volume Histogram \\
EUD & Equivalent Uniform Dose \\
FU & Fraction of Uptake \\
MIRD & Medical Internal Radiation Dose \\
NTCP & Normal Tissue Complication Probability \\
SF & Shunt Fraction \\
T/N & Tumor-to-Normal ratio \\
TCP & Tumor Control Probability \\
\bottomrule
\end{tabular}

\subsection{Propiedades del Isotopo}

El $^{90}$Y es un emisor $\beta^-$ puro con las siguientes caracteristicas:

\begin{table}[H]
\centering
\begin{tabular}{l c l}
\toprule
\textbf{Propiedad} & \textbf{Valor} & \textbf{Unidad} \\
\midrule
Vida media $T_{1/2}$ & 64.1 & horas \\
Constante de desintegracion $\lambda$ & $3.00 \times 10^{-6}$ & s$^{-1}$ \\
Vida media $\tau$ & 332,916 & s \\
Energia maxima $\beta^-$ $E_{\max}$ & 2.28 & MeV \\
Energia media $\beta^-$ $\bar{E}_\beta$ & 0.9336 & MeV \\
Alcance maximo en tejido & $\sim 11$ & mm \\
Alcance medio & $\sim 2.5$ & mm \\
\bottomrule
\end{tabular}
\end{table}

\subsection{Dosis Absorbida}

La dosis absorbida $D$ se define como la energia depositada por unidad de masa:

\begin{equation}\label{eq:dosis}
    D = \frac{dE}{dm} \quad [\text{J/kg} = \text{Gy}]
\end{equation}

Para $^{90}$Y en dosimetria interna, la dosis en un voxel depende de la actividad acumulada,
la energia liberada por desintegracion ($\bar{E}_\beta$), la fraccion de energia absorbida
localmente (esencialmente 1 para $\beta^-$ de $^{90}$Y) y la masa del voxel.

\subsection{Modelo MIRD (Partition Model)}

El modelo MIRD de particion provee estimaciones analiticas rapidas:

\begin{align}
    T/N &= \frac{\text{actividad\_PET\_tumor}}{\text{actividad\_PET\_higado}} \\
    \text{FU}_{\text{normal}} &= (1 - \text{SF})\,
          \frac{V_{\text{higado}}}{T/N \cdot V_{\text{tumor}} + V_{\text{higado}}} \\
    \text{FU}_{\text{tumor}} &= (1 - \text{SF})\,
          \frac{T/N \cdot V_{\text{tumor}}}{T/N \cdot V_{\text{tumor}} + V_{\text{higado}}} \\
    D_{\text{higado}} &= \frac{A \cdot k \cdot \text{FU}_{\text{normal}}}{m_{\text{higado}}} \\
    D_{\text{tumor}} &= \frac{A \cdot k \cdot \text{FU}_{\text{tumor}}}{m_{\text{tumor}}}
\end{align}

\subsection{Variables del Modelo MIRD --- Tabla Completa}

\begin{table}[H]
\centering
\begin{tabular}{lll}
\toprule
\textbf{Variable} & \textbf{Descripcion} & \textbf{Unidad} \\
\midrule
$T/N$ & Tumor-to-Normal ratio (razon de captacion) & adimensional \\
SF & Shunt Fraction (fraccion que escapa al pulmon) & adimensional \\
$\text{FU}_{\text{normal}}$ & Fraction of Uptake en higado sano & adimensional \\
$\text{FU}_{\text{tumor}}$ & Fraction of Uptake en tumor & adimensional \\
$V_{\text{higado}}$ & Volumen del higado & cm$^3$ \\
$V_{\text{tumor}}$ & Volumen del tumor & cm$^3$ \\
$m_{\text{higado}}$ & Masa del higado & g \\
$m_{\text{tumor}}$ & Masa del tumor & g \\
$A$ & Actividad administrada & GBq \\
$k$ & Constante MIRD para $^{90}$Y ($= 48.98$) & J/s \\
$D_{\text{higado}}$ & Dosis absorbida en higado sano & Gy \\
$D_{\text{tumor}}$ & Dosis absorbida en tumor & Gy \\
$D$ & Dosis absorbida generica & Gy \\
$\bar{E}_\beta$ & Energia media de la particula beta & MeV \\
$\lambda$ & Constante de desintegracion & s$^{-1}$ \\
$\tau$ & Vida media del isotopo & s \\
\bottomrule
\end{tabular}
\end{table}

\subsection{Constantes Fisicas}

\begin{align}
    \lambda &= \frac{\ln 2}{T_{1/2}} \approx 3.00 \times 10^{-6} \text{ s}^{-1} \\
    \tau &= \frac{1}{\lambda} = \frac{T_{1/2}}{\ln 2} \approx 332\,916 \text{ s}
\end{align}

\newpage

% ============================================================
% REFERENCIAS
% ============================================================
\section{Referencias}

\begin{enumerate}
    \item Wasserthal, J., et al. ``TotalSegmentator: Robust Segmentation
          of 104 Anatomic Structures in CT Images.'' \textit{Radiology
          Artificial Intelligence}, 2023.
    \item ICRP Publication 110. ``Adult Reference Computational
          Phantoms.'' \textit{Annals of the ICRP}, 2009.
    \item ICRU Report 44. ``Tissue Substitutes in Radiation Dosimetry
          and Measurement.'' 1989.
    \item Klein, S., et al. ``Elastix: a toolbox for intensity-based
          medical image registration.'' \textit{IEEE Trans. Med.
          Imaging}, 2010.
    \item Fedorov, A., et al. ``3D Slicer as an image computing
          platform for the Quantitative Imaging Network.''
          \textit{Magn. Reson. Imaging}, 2012.
    \item Gulec, S.A., et al. ``Hepatic radiation dosimetry in
          $^{90}$Y-microsphere treatment.'' \textit{Med. Phys.}, 2003.
    \item Chiesa, C., et al. ``The dosimetric journey of $^{90}$Y
          radioembolization: from MIRD to voxel dosimetry.''
          \textit{EJNMMI Physics}, 2017.
    \item Documentacion del proyecto 3Dosim: pipeline\_mod1.py,
          pet\_dicom\_reader.py, couch\_remover.py,
          pet\_registration.py, fusion\_dialog.py, segmentation.py,
          tumor\_creator.py, validation.py, labelmap\_exporter.py,
          checkpoint.py, ai\_supervisor.py.
\end{enumerate}

\end{document}
